\section{What standard DFPT actually assumes}
\label{sec:dfpt}

\subsection{The perturbation and the matrix element}

Let $u_{\nu\vq}$ be the normal coordinate of phonon branch $\nu$ and wavevector
$\vq$.  After including the conventional phonon normalization factor, define the
self-consistent Kohn--Sham perturbation operator
\begin{equation}
  D^{\ks}_{\nu\vq}
  \equiv \partial_{\nu\vq} V_{\ks}
  = \partial_{\nu\vq} V_{\mathrm{ion}}
  + \partial_{\nu\vq} V_{\mathrm{Hxc}} .
  \label{eq:dks}
\end{equation}
The DFPT electron--phonon matrix element is
\begin{equation}
  g^{\dfpt}_{mn\nu}(\vk,\vq)
  = \langle \psi_{m,\vk+\vq} |
    D^{\ks}_{\nu\vq}
    | \psi_{n,\vk} \rangle .
  \label{eq:gdfpt}
\end{equation}
The perturbation is a matrix in band, orbital, spin, and reciprocal-lattice
indices.  Therefore standard DFPT does not assume that every electron sees the
same scalar shift.

The first-order orbital response is obtained from a Sternheimer equation of the
schematic form
\begin{equation}
  \bigl(\hat H_{\ks}-\varepsilon_{n\vk}\bigr)
  |\partial_{\nu\vq}\psi_{n\vk}\rangle
  = -\hat Q_{\vk+\vq}
  \bigl(\partial_{\nu\vq}V_{\ks}
  -\partial_{\nu\vq}\varepsilon_{n\vk}\bigr)
  |\psi_{n\vk}\rangle,
  \label{eq:sternheimer}
\end{equation}
where $\hat Q$ projects outside the occupied state or selected variational
subspace.  The induced density updates $V_{\mathrm{Hxc}}$ until self-consistency.

\subsection{Static electronic equilibrium}

In ordinary adiabatic DFPT, the electrons are assumed to remain in the
instantaneous ground state associated with a slowly displaced lattice
\citep{Gonze1995,Baroni2001}.  Equivalently, the electronic response entering
\cref{eq:sternheimer} is evaluated at zero external frequency.  The electronic
denominator contains Kohn--Sham energy differences, but not a phonon-frequency
shift such as $\varepsilon_{m,\vk+\vq}-\varepsilon_{n\vk}\pm\omega_{\nu\vq}$.

This assumption is usually sensible when electronic relaxation and interband
energy scales are fast compared with the phonon period.  It can fail for low
carrier densities, narrow bands, small Fermi energies, polar optical modes,
collective modes, or any situation in which $\omega_{\nu\vq}$ is not negligible
on the scale governing the electronic response.

\subsection{The response equation behind DFPT screening}

For a scalar density perturbation, suppressing matrix indices for clarity,
\begin{align}
  \delta V_{\ks}
  &= \delta V_{\mathrm{ion}}+(v+f_{\mathrm{xc}})\delta n,
  \label{eq:vksresponse}\\
  \delta n &= \chi_s\,\delta V_{\ks},
  \label{eq:chisresponse}
\end{align}
where $v$ is the Coulomb interaction, $f_{\mathrm{xc}}$ is the static
exchange--correlation kernel, and $\chi_s$ is the independent Kohn--Sham density
response.  Solving the two equations gives
\begin{equation}
  \delta V_{\ks}
  = \bigl[1-(v+f_{\mathrm{xc}})\chi_s\bigr]^{-1}
    \delta V_{\mathrm{ion}},
  \qquad
  g^{\dfpt}
  = \frac{g^0}{1-(v+f_{\mathrm{xc}})\chi_s}
  \label{eq:gdfptscalar}
\end{equation}
in a translationally invariant scalar problem.  In a crystal these quantities are
matrices in reciprocal lattice vectors and may also carry spin or sublattice
indices.

\subsection{The precise missing term}

Choose the convention
\begin{equation}
  G^{-1}(\omega)=
  \omega+\mu-H_{\ks}-\Sigma^{\mathrm{corr}}(\omega),
  \label{eq:gconvention}
\end{equation}
where $\Sigma^{\mathrm{corr}}$ is the self-energy correction beyond the chosen
Kohn--Sham reference, with the corresponding double counting removed.  The
many-body perturbation of the quasiparticle equation contains
\begin{equation}
  D^{\mb}_{\nu\vq}(\omega)
  = \partial_{\nu\vq}
    \bigl[V_{\ks}+\Sigma^{\mathrm{corr}}(\omega)\bigr].
  \label{eq:dmb}
\end{equation}
Thus the central vertex term absent from standard DFPT is
$\partial_{\nu\vq}\Sigma^{\mathrm{corr}}$, together with its frequency
dependence and consistent wavefunction renormalization.  A large equilibrium
self-energy does not logically imply a large derivative with respect to a
particular phonon.  This distinction is the first reason why the label ``strongly
correlated'' alone cannot predict the DFPT error.  A consistent separation between
the DFT starting point and many-body electron--phonon corrections is discussed in
detail by \citet{Marini2015}.

\begin{keypoint}[Operational limitation of DFPT]
Standard DFPT is a static Kohn--Sham response.  Its reliability for a physical
quasiparticle vertex depends on whether the phonon derivative of the beyond-DFT
self-energy is small, already represented by the chosen static density
functional, or cancelled by a constraint in the relevant operator channel.
\end{keypoint}
